% RecSys Odyssey repository-backed lecture note.
% Add figures with: \coursefigure{figures/name.svg}{Caption}
\section{Blend sources and inspect misses}

\begin{abstract}
One retriever rarely covers every intent, item age and failure mode.
\end{abstract}

\subsection{Concept}
Production candidate sets often merge several sources: two-tower similarity, item-to-item neighbors, recent popularity, subscriptions, editorial pools and exploration. Source quotas prevent one retriever from crowding out all others; deduplication and eligibility checks create a clean union. Debugging should trace each final item back to its source and inspect relevant items that were missed. Retrieval is complete only when its outputs are observable and refreshable.

\begin{equation*}
C = dedupe(C_two-tower \cup C_item-item \cup C_popular \cup C_fresh \cup C_explore)
\end{equation*}

\subsection{Key terms}
\begin{itemize}
\item \textbf{Candidate source.} One retrieval strategy contributing items to the merged shortlist.
\item \textbf{Source quota.} A limit or reservation controlling each source’s contribution.
\item \textbf{Lineage.} A trace of where an item entered and how it moved through the pipeline.
\end{itemize}
